\begin{initiativkarte}{IAESTE / AIESEC – Internationale Praktika}

  % Bild (optional)
  \IfFileExists{bilder/iaeste.jpg}{%
    \begin{center}
      \includegraphics[width=\linewidth,height=5cm,keepaspectratio]{bilder/iaeste.jpg}
    \end{center}
    \vspace{0.4cm}
  }{}

  % Kurzbeschreibung
  \textit{\textcolor{hawgray}{IAESTE vermittelt internationale Fachpraktika in Ingenieur- und Naturwissenschaften in über 80 Ländern weltweit.}}

  \vspace{0.5cm}
  \hawrule

  % Beschreibung
  \textbf{\textcolor{hawblue}{Was ist IAESTE / AIESEC?}}\\[0.3cm]

  IAESTE ist ein internationales, gemeinnütziges Praktikantenaustauschprogramm für Studierende der Ingenieur- und Naturwissenschaften.

  Es vermittelt bezahlte Fachpraktika im Ausland und ermöglicht dadurch wertvolle internationale Berufserfahrung.

  Typische Praktika dauern zwischen 2 und 3 Monaten, sind aber auch im Zeitraum von 4 Wochen bis zu 12 Monaten möglich.

  \vspace{0.3cm}

  \textbf{Leistungen von IAESTE:}
  \begin{itemize}
      \item Vermittlung kostenloser internationaler Praktika
      \item Praktikumsplätze in über 80 Ländern
      \item Unterstützung bei Bewerbung, Visa und Organisation
      \item Betreuung durch Lokalkomitees im Heimat- und Zielland
      \item Hilfe bei Wohnungssuche und Integration im Ausland
  \end{itemize}

  \vspace{0.3cm}

  Die Lebenshaltungskosten werden in der Regel durch die Praktikumsvergütung gedeckt. IAESTE selbst zahlt keine Zuschüsse.

  Teilweise sind Fahrtkostenzuschüsse über den DAAD möglich.

  \vspace{0.5cm}

  \textbf{Voraussetzungen:}
  \begin{itemize}
      \item mindestens 3. Semester
      \item erste praktische Erfahrung im Fachbereich empfohlen
      \item gute Englischkenntnisse
      \item Altersgrenze: 30 Jahre
  \end{itemize}

  \vspace{0.5cm}

  % Tags
  \tag{International}
  \tag{Praktikum}
  \tag{Karriere}
  \tag{Engineering}
  \tag{Austausch}

  \vspace{0.6cm}
  \hawrule

  % Infozeilen
  \infozeile{\faUsers}{Zielgruppe: Studierende der Ingenieur- und Naturwissenschaften}
  \infozeile{\faGraduationCap}{Semester: Ab ca. 3. Semester}
  \infozeile{\faMapMarker*}{Ort: TU Hamburg (IAESTE Lokalkomitee Hamburg)}
  \infozeile{\faGlobe}{Website: \href{https://www.iaeste-hamburg.de}{iaeste-hamburg.de}}

\end{initiativkarte}